% !TeX root = ../../Book1.tex
\chapter{\texorpdfstring{\(L^p\)}{Lᵖ} 空间}
\label{b1:ch:10}
\BookChapterNumber{b1:ch:10}
% 资料组：Axler2020Measure 7A--B、9.42；Fremlin2016Measure 246--247；
% BallMurat1989Biting为正式文献定位，未声称实际读到受限全文。
% 证明义务：等价类与零集、p=∞端点、可求和误差的完备性、稠密性的测度范围、
% sigma有限对偶表示的局部RN与密度范数估计、L1中集中/逃逸、Chacon尾部质量构造。
\begin{chapterabstract}
测度论把函数按几乎处处相等归为同一对象，泛函分析则要求为这些对象选定范数和连续测试。
而 \(L^p\) 空间把两种结构连接起来：指数 \(p\) 决定大值受到多强的惩罚，
也决定函数可以与哪些测试函数相乘积分。本章先建立范数、完备性和稠密逼近，
再用 Radon--Nikodym 定理识别对偶空间，把前一章的自反性与弱紧性落实到函数上。
两个端点不能仅靠形式代入处理：\(L^1\) 容许质量集中；在典型的无限维情形中，\(L^\infty\) 的连续对偶通常严格大于由 \(L^1\) 密度产生的积分泛函。
选读部分以一致可积性和 Chacon 引理进一步说明这些边界现象。
\end{chapterabstract}

% 本章目标（不计入编号）
\begin{learninggoals}
\begin{itemize}
\item\label{b1:item:10-goal-norm}在几乎处处等价类上使用 \(L^p\) 范数，掌握 Hölder、Minkowski 不等式和完备性的证明。
\item\label{b1:item:10-goal-density}根据指数和测度条件选择简单函数、截断函数或连续紧支集函数作逼近，不混淆有限 \(p\) 与无穷端点。
\item\label{b1:item:10-goal-duality}用积分配对识别对偶，并理解 Radon--Nikodym 密度为何具有恰好的可积指数。
\item\label{b1:item:10-goal-weak}在 \(1<p<\infty\) 的范围内使用自反性抽取弱子列，识别能量与约束在取极限时的变化。
\item\label{b1:item:10-goal-endpoints}区别 \(L^1\) 中的集中、无穷远处的质量逃逸和普通振荡，理解一致可积性如何补充范数有界性。
\end{itemize}
\end{learninggoals}

前面已经多次使用 \(L^1\) 和 \(L^2\)。本章不是把积分重新定义一遍，
而是统一处理整个尺度，并把“积分有限”提升为可以求极限、作逼近和识别连续泛函的空间结构。
基础部分允许一般测度空间；对偶表示和自反性的主线证明明确在 \(\sigma\)-有限测度下展开。
选读的一致可积性紧性判别先处理有限测度空间，随后说明非有限空间还缺少哪一种控制。
每次使用结论时，应保留这些测度假设，不能只记住指数范围。

第一遍宜把\secref{b1:sec:1001}、\secref{b1:sec:1002}连读，
再将\secref{b1:sec:1003}的表示证明与\chapref{b1:ch:06}对照：可数可加性负责产生测度，
绝对连续性负责产生密度，最后的测试函数负责确定密度的范数。
\secref{b1:sec:1004}把这些结果代入上一章的抽象工具。
若当前目的只是准备 \(p>1\) 的 Sobolev 空间，可以先略过末节；
处理仅有一次幂积分估计的近似解时，再回来阅读端点紧性。

Sheldon Axler 的《Measure, Integration \& Real Analysis》%
\cite[Chapter 7; 9.42 dual space of \(L^p(\mu)\) is \(L^{p'}(\mu)\)]{Axler2020Measure}可作为这一章的主线伴读。
其中等价类、范数不等式与一般对偶表示分处不同位置，适合用来复查测度论与空间结构的衔接。
一致可积性和 \(L^1\) 紧性的较一般形式，参阅 Fremlin 的《Measure Theory》%
\cite[Volume 2, \S\S246--247]{Fremlin2016Measure}。
中文术语和对偶理论亦可配合许全华等《泛函分析讲义》%
\cite[第9.5节]{Xu2017Fanhan}阅读；此处仅作已核定目录范围内的伴读推荐。

% !TeX root = ../../Book1.tex
\section{\texorpdfstring{\(L^p\)}{Lᵖ} 的基本结构}
\label{b1:sec:1001}

% TODO: 撰写本节正文。

\subsection{等价类}
\label{b1:subsec:100101}

待填充。

\subsection{\texorpdfstring{\(L^p\)}{Lp} 范数}
\label{b1:subsec:100102}

待填充。

\subsection{Hölder 不等式}
\label{b1:subsec:100103}

待填充。

\subsection{Minkowski 不等式}
\label{b1:subsec:100104}

待填充。

\subsection{完备性}
\label{b1:subsec:100105}

待填充。

% !TeX root = ../../Book1.tex
\section{稠密性}
\label{b1:sec:1002}

% 来源：Axler2020Measure，7A习题17--19，第200--201页，结合本卷第4章简单逼近与第5章DCT。
% 改编：把实直线上的连续紧支集逼近推广至Rn，用第3章正则性和距离函数完整证明；
% 区分有限p与无穷端点，有限测度非零区域由函数自身选取，不额外假设sigma有限。
完备性保证逼近过程有极限，稠密性则说明可以从哪些容易处理的函数开始逼近。这里依次减少函数的复杂性：先使值域有限，再控制幅值与非零区域，最后在欧氏空间中要求连续性。每次增加限制，都要检查误差仍能在所选的 \(L^p\) 范数下趋于零；有限指数与无穷端点在这一点上有显著差别。

\subsection{简单函数}
\label{b1:subsec:100201}

\begin{proposition}\label{b1:prop:lp-simple-dense}
对任意测度空间及 \(1\leq p\leq\infty\)，属于 \(L^p(\mu)\) 的简单函数在 \(L^p(\mu)\) 中稠密。当 \(p<\infty\) 时，这些简单函数还可以取为在某个有限测度集合外恒为零。
\end{proposition}
\begin{proof}
先设 \(p<\infty\)、\(f\geq0\)。由\cref{b1:thm:nonnegative-simple-approximation}，存在非负简单函数 \(s_n\uparrow f\)。因
\[
 |f-s_n|^p\longrightarrow0,\quad |f-s_n|^p\leq f^p\in L^1,
\]
\cref{b1:thm:dominated-convergence}给出 \(\|f-s_n\|_p\to0\)。写出 \(s_n\) 的非零规范表示 \(s_n=\sum_{j=1}^ma_j\mathbf1_{A_j}\)，其中 \(a_j>0\)。由
\[
 a_j^p\mu(A_j)\leq\int_Xs_n^p\,\mathrm d\mu
 \leq\int_Xf^p\,\mathrm d\mu<\infty
\]
知每个 \(A_j\) 都有有限测度，其有限并容纳 \(s_n\) 的非零区域。对实函数分别逼近正、负部，对复函数再分别处理实、虚部；各部分的 \(p\) 次幂均受 \(|f|^p\) 控制，Minkowski 不等式保证组合后的误差趋零。有限个非零区域的并仍有有限测度。

若 \(p=\infty\)，先选取处处有界的可测代表：把 \(\{|f|>\|f\|_\infty\}\) 上的值改成零即可。对实、虚部各自按长度 \(2^{-n}\) 的区间向下取整，得到有限值的可测函数 \(s_n\)，并有
\[
 |f-s_n|\leq\sqrt2\,2^{-n}
\]
处处成立；实情形还可去掉 \(\sqrt2\)。所以 \(\|f-s_n\|_\infty\to0\)。这里并未要求 \(s_n\) 的非零区域具有有限测度。
\end{proof}

“简单”限制的是值域中数值的个数，而不是函数非零处的测度。比如 \(\mathbb R^d\) 上的常函数 \(1\) 已经是简单函数，也属于 \(L^\infty\)，却不属于任何有限指数的 \(L^p\)。因此在无限测度空间上，不能把任意简单函数都当作有限 \(p\) 的逼近候选。对有限测度集 \(A,B\)，一个常用的精确公式是
\[
 \|\mathbf1_A-\mathbf1_B\|_p^p
 =\mu(A\mathbin{\triangle}B)\quad(1\leq p<\infty).
\]
这把集合逼近的测度误差，直接转换成函数逼近的范数误差。

\subsection{截断}
\label{b1:subsec:100202}

对实值函数，最熟悉的截断是把过大的正、负值分别压到 \(M\) 与 \(-M\)。为了兼顾复值函数，我们沿每个函数值的方向截短其模，定义
\[
 T_Mf(x)=
 \begin{cases}
 f(x),&|f(x)|\leq M,\\
 M\dfrac{f(x)}{|f(x)|},&|f(x)|>M,
 \end{cases}
 \quad M>0.
\]
这是可测函数，且 \(|T_Mf|\leq|f|\)、\(|T_Mf|\leq M\)。

\begin{proposition}\label{b1:prop:lp-truncation}
若 \(1\leq p<\infty\) 且 \(f\in L^p(\mu)\)，则 \(T_Mf\to f\) 于 \(L^p\)。更一般地，若 \(E_n\in\mathcal A\) 递增，且 \(f=0\) 几乎处处于 \(X\setminus\bigcup_nE_n\)，则 \(f\mathbf1_{E_n}\to f\) 于 \(L^p\)。
\end{proposition}
\begin{proof}
第一项的误差满足
\[
 |f-T_Mf|^p=(|f|-M)_+^p\longrightarrow0,
 \quad |f-T_Mf|^p\leq|f|^p.
\]
第二项的误差是 \(|f|^p\mathbf1_{X\setminus E_n}\)，也几乎处处趋零并受 \(|f|^p\) 控制。两次使用\cref{b1:thm:dominated-convergence}即可。
\end{proof}

例如在 \((0,1)\) 上，\(f(x)=x^{-\alpha}\) 属于 \(L^p\) 当且仅当 \(\alpha p<1\)（\(\alpha>0\)）。满足这一条件时，有界截断可以在 \(L^p\) 范数下逼近这个无界函数，尽管每次截断留下的逐点误差仍没有有限上界。积分范数允许极大的误差落在越来越小的区域中，控制收敛正是对这种容许方式的定量依据。

\(L^\infty\) 中，幅值截断对给定本性有界函数最终等于该函数，因而没有困难；困难在于区域截断。即使 \(E_n\uparrow X\)，也未必有 \(\|f-f\mathbf1_{E_n}\|_\infty\to0\)。在 \(\mathbb R\) 上取 \(f=1\)、\(E_n=[-n,n]\)，误差范数始终为 \(1\)。集合尾部虽然在逐点意义下消失，其上保留的误差高度并没有降低。

\subsection{有限测度支集函数}
\label{b1:subsec:100203}

一般可测空间尚无拓扑，因而本小节所说的“有限测度支集函数”，仅指存在 \(E\in\mathcal A\)、\(\mu(E)<\infty\)，使函数在 \(X\setminus E\) 上几乎处处为零。这里不取非零集合的拓扑闭包，也不要求空间本身能被可数个有限测度集覆盖。

\begin{proposition}\label{b1:prop:lp-finite-support-dense}
若 \(1\leq p<\infty\)，则有界且在某个有限测度集合外为零的可测函数在 \(L^p(\mu)\) 中稠密。对每个给定的可测代表 \(f\in L^p\)，可具体取
\[
 E_n=\{\,1/n<|f|\leq n\,\},\quad f_n=f\mathbf1_{E_n}.
\]
\end{proposition}
\begin{proof}
由水平集估计，
\[
 \mu(E_n)\leq\mu(\{|f|>1/n\})
 \leq n^p\int_X|f|^p\,\mathrm d\mu<\infty.
\]
又有 \(|f_n|\leq n\)，且 \(E_n\uparrow\{f\ne0\}\)，所以 \(f_n\to f\) 处处。最后，\(|f-f_n|^p\leq|f|^p\)，控制收敛给出范数逼近。
\end{proof}

这一构造由函数决定局部化区域。即使 \(X\) 的测度非常大，每个 \(L^p\) 函数的非零区域仍包含在可数个有限测度集的并中。不同函数可能需要不同的区域，不能因此反推 \(\mu\) 是 \(\sigma\)-有限测度。例如在不可数集合的计数测度下，每个有限指数的 \(L^p\) 函数至多有可数多个非零值，但整个空间不是 \(\sigma\)-有限的。

有限测度支集也不等于拓扑紧支集。在 \(\mathbb R\) 中，可取分布在各整数附近、总长度有限而整体无界的开区间并；其示性函数属于每个有限指数的 \(L^p\)，非零区域却不在任何紧集内。欧氏空间中的紧支集逼近还要进一步使用测度与拓扑之间的正则性。

上述稠密性不能推广到 \(L^\infty\)。若 \(g\) 在某个有限 Lebesgue 测度集外为零，则 \(\|1-g\|_\infty\geq1\)，因为该集的补集具有正测度。这不仅否定某一个截断序列，而且否定用全部有限测度支集函数去逼近常函数 \(1\) 的可能性。

\subsection{欧氏空间中的光滑逼近预告}
\label{b1:subsec:100204}

现在固定 \(\mathbb R^d\) 上的 Lebesgue 测度。连续函数 \(u\) 的\term[zhiji]{支集}[support]定义为
\[
 \operatorname{supp}u=\overline{\{\,x\in\mathbb R^d\mid u(x)\ne0\,\}}.
\]
记 \(C_c(\mathbb R^d)\) 为具有紧支集的连续实值或复值函数空间。紧集有有限测度，而连续函数在紧集上有界，所以 \(C_c\subset L^p\) 对每个 \(1\leq p\leq\infty\) 都成立。
\symidx{\(\operatorname{supp}u\)：连续函数的支集}
\symidx{\(C_c(\mathbb R^d)\)：欧氏空间上的连续紧支集函数空间}

\begin{proposition}\label{b1:prop:lp-continuous-dense}
若 \(1\leq p<\infty\)，则 \(C_c(\mathbb R^d)\) 在 \(L^p(\mathbb R^d)\) 中稠密。
\end{proposition}
\begin{proof}
由\cref{b1:prop:lp-simple-dense}和 Minkowski 不等式，只须把有限测度集 \(A\) 的示性函数逼近为连续紧支集函数。给定 \(\eta>0\)，由\cref{b1:thm:inner-regular}取紧集 \(K\subset A\)，使 \(m(A\setminus K)<\eta\)。若 \(K=\varnothing\)，零函数已经满足所需误差估计。否则，由外正则性取包含 \(K\) 的开集，再与一个包含 \(K\) 的大开球相交，可得到有界开集 \(U\supset K\)，满足 \(m(U\setminus K)<\eta\)。

记 \(d(x,F)=\inf_{y\in F}|x-y|\)。到非空闭集的距离函数连续，因为 \(|d(x,F)-d(y,F)|\leq|x-y|\)。定义
\[
 u(x)=\frac{d(x,U^c)}{d(x,K)+d(x,U^c)}.
\]
分母处处为正：若两个距离同时为零，闭性便给出 \(x\in K\cap U^c\)，与 \(K\subset U\) 矛盾。故 \(u\) 连续，\(0\leq u\leq1\)，在 \(K\) 上等于 \(1\)，在 \(U^c\) 上等于 \(0\)，并且 \(\operatorname{supp}u\subset\overline U\) 为紧集。函数 \(u\) 与 \(\mathbf1_A\) 只有在 \((A\setminus K)\cup(U\setminus K)\) 上可能不同，差的模不超过 \(1\)，所以
\[
 \|\mathbf1_A-u\|_p^p
 \leq m(A\setminus K)+m(U\setminus K)<2\eta.
\]
让 \(\eta\) 足够小便可达到任意给定误差。对有限和 \(s=\sum_{j=1}^ra_j\mathbf1_{A_j}\)，分别选择这些连续逼近，使 \(\sum_j|a_j|\cdot\|\mathbf1_{A_j}-u_j\|_p\) 任意小。有限和 \(\sum_ja_ju_j\) 仍连续且具有紧支集，Minkowski 不等式完成证明。
\end{proof}

这一证明把正则性用于跳跃发生的边界：连续函数不必逐点复制示性函数，只须把过渡安排在测度很小的区域中。它建立了从可测函数到连续函数的桥，但还没有建立从连续函数到光滑函数的桥。以后引入卷积磨光时，我们将构造具有紧支集的无穷次可微逼近，并证明磨光误差在 \(L^p\) 范数下趋零；本节不把尚未证明的磨光性质用作稠密性依据。

无穷端点仍然有障碍，而且即使没有远处尾部，跳跃也足以阻止逼近。在 \(\mathbb R\) 上令 \(f=\mathbf1_{(0,1)}\)。对任何连续函数 \(g\)，都有 \(\|f-g\|_\infty\geq1/2\)。否则可取 \(a<1/2\)，使 \(|f-g|\leq a\) 几乎处处。于 \(0\) 左侧取趋于 \(0\) 的非例外点，并用连续性，得到 \(|g(0)|\leq a\)；于右侧同样得到 \(|1-g(0)|\leq a\)。三角不等式便给出 \(1\leq2a<1\)，矛盾。因而连续紧支集函数乃至光滑紧支集函数都不可能在整个 \(L^\infty\) 中稠密。

% !TeX root = ../../Book1.tex
\section{\texorpdfstring{\(L^p\)}{Lᵖ} 对偶}
\label{b1:sec:1003}

Hölder 不等式已经说明，一个可积乘积能够产生连续线性泛函。对偶表示要回答反方向的问题：每个连续线性泛函是否都来自这样的积分？本节以 \(\sigma\)-有限正测度为主线，用第\ref{b1:ch:06}章的 Radon--Nikodym 定理构造表示密度。有限测度块上的构造可参看 Axler 的《Measure, Integration \& Real Analysis》\cite[第 7B 节及定理 9.42]{Axler2020Measure}。该书还处理了 \(1<p<\infty\) 时更一般的测度；以下只证明所明列的 \(\sigma\)-有限版本，不把本节的证明范围误当成定理成立的必要条件。

% 实读来源：Axler2020Measure，7.25--7.26，印刷 pp.206--207；
% 一般Lp的RN证明不在第7章，而在9.42，有限测度部分为 pp.275--276。
% 本节改写为sigma有限块拼接，另完整证明p=1端点与复数相位测试。

\subsection{\texorpdfstring{\(1<p<\infty\)}{1<p<∞}}
\label{b1:subsec:100301}

固定标量域 \(\mathbb K=\mathbb R\) 或 \(\mathbb C\)。设 \(1<p<\infty\)，记 \(q=p'=p/(p-1)\)。给定 \(g\in L^q(\mu;\mathbb K)\)，定义
\[
 \Lambda_g:L^p(\mu;\mathbb K)\longrightarrow\mathbb K,
 \quad \Lambda_g(f)=\int_Xfg\,\mathrm d\mu.
\]
这里不在 \(g\) 上加共轭，因此 \(f\mapsto\Lambda_g(f)\) 和 \(g\mapsto\Lambda_g\) 都线性。若采用 Hilbert 空间的写法 \(\int f\overline h\,\mathrm d\mu\)，对应的密度就是 \(g=\overline h\)，两种表示的数学内容一致，但参数到泛函的线性性质不同。

\begin{proposition}\label{b1:prop:lp-dual-isometry}
设 \(\mu\) 为任意正测度，\(1<p<\infty\)。上述映射
\[
 D_p:L^q(\mu)\longrightarrow (L^p(\mu))^*,
 \quad D_pg=\Lambda_g
\]
是线性等距嵌入。
\end{proposition}
\begin{proof}
由\cref{b1:thm:lp-holder}，乘积属于 \(L^1(\mu)\)，且
\[
 |\Lambda_g(f)|\le\|g\|_q\|f\|_p.
\]
改变任一个函数的几乎处处代表不改变积分，所以泛函在等价类上良定义。

为得到反向估计，令
\[
 \vartheta_g(x)=
 \begin{cases}
  \overline{g(x)}/|g(x)|,&g(x)\ne0,\\
  0,&g(x)=0.
 \end{cases}
\]
实情形中的共轭不起作用，而两种情形都满足 \(\vartheta_gg=|g|\)。若 \(\|g\|_q>0\)，取
\[
 f_0=\dfrac{\vartheta_g|g|^{q-1}}{\|g\|_q^{q-1}}.
\]
因为 \((q-1)p=q\)，有 \(\|f_0\|_p=1\) 及
\(\Lambda_g(f_0)=\|g\|_q\)。故 \(\|\Lambda_g\|=\|g\|_q\)；\(g=0\) 的情形直接成立。积分的线性性给出 \(D_p\) 线性，等距性则保证单射。
\end{proof}

这一步不要求 \(\sigma\)-有限性，而且给出了达到对偶范数的测试函数。尚未解决的是满射性：一个抽象的连续泛函并没有预先给定密度。

\begin{theorem}[\(L^p\) 对偶表示]\label{b1:thm:lp-dual}
设 \((X,\mathcal A,\mu)\) 为 \(\sigma\)-有限正测度空间，\(1<p<\infty\)，\(q=p'\)。则每个 \(\Lambda\in(L^p(\mu;\mathbb K))^*\) 都唯一表示为
\[
 \Lambda(f)=\int_Xfg\,\mathrm d\mu
 \quad(f\in L^p(\mu;\mathbb K)),
\]
其中 \(g\in L^q(\mu;\mathbb K)\)，且 \(\|\Lambda\|=\|g\|_q\)。
\end{theorem}

唯一性及等范数已由前一命题给出。下面用可测集合上的取值制造密度，再从泛函的有界性中读出密度的可积指数。

\subsection{Radon–Nikodym 方法}
\label{b1:subsec:100302}

若 \(\mu(X)=\infty\)，一般不能写 \(\Lambda(\mathbf1_X)\)，因为 \(\mathbf1_X\) 可能不在 \(L^p\) 中。因此构造首先在有限测度空间上进行；这个局部步骤也将用于 \(p=1\)。

\begin{lemma}\label{b1:lem:lp-functional-local-density}
设 \(\mu(X)<\infty\)，\(1\le p<\infty\)，\(\Lambda\in(L^p(\mu;\mathbb K))^*\)。则存在 \(g\in L^1(\mu;\mathbb K)\)，使
\[
 \Lambda(v)=\int_Xvg\,\mathrm d\mu
\]
对每个有界可测函数 \(v\) 成立。
\end{lemma}
\begin{proof}
对 \(E\in\mathcal A\) 定义 \(\nu(E)=\Lambda(\mathbf1_E)\)。若各 \(E_j\) 两两不交，则
\[
 \left\|\mathbf1_{\bigcup_jE_j}-\sum_{j=1}^n\mathbf1_{E_j}\right\|_p^p
 =\mu\left(\bigcup_{j>n}E_j\right)\longrightarrow0.
\]
最后一步使用了有限测度性。因此 \(\Lambda\) 的连续性给出
\(\nu(\bigcup_jE_j)=\sum_j\nu(E_j)\)。又当 \(\mu(E)=0\) 时，\(\mathbf1_E\) 是零等价类，故 \(\nu(E)=0\)。于是 \(\nu\) 是绝对连续于 \(\mu\) 的符号测度或复测度。

实情形由\cref{b1:cor:signed-radon-nikodym}得到 \(g\in L^1(\mu;\mathbb R)\)。复情形分别对 \(\operatorname{Re}\nu\)、\(\operatorname{Im}\nu\) 使用同一结果，得到实密度 \(a,b\in L^1(\mu)\)，令 \(g=a+\mathrm i b\)。两种情形都满足
\[
 \Lambda(\mathbf1_E)=\int_Eg\,\mathrm d\mu.
\]
有限线性组合将等式推广到简单函数。对有界可测 \(v\)，取一致逼近它的简单函数 \(s_n\)。由于
\[
 \|s_n-v\|_p\le\mu(X)^{1/p}\|s_n-v\|_\infty,
 \quad
 \left|\int_X(s_n-v)g\,\mathrm d\mu\right|
 \le\|s_n-v\|_\infty\|g\|_1,
\]
两边均可取极限，得到结论。
\end{proof}

\begin{proof}[定理\ref{b1:thm:lp-dual}的证明]
取两两不交的可测集 \(A_j\)，使 \(X=\bigcup_jA_j\)、\(\mu(A_j)<\infty\)，并令 \(E_N=\bigcup_{j=1}^NA_j\)。把 \(L^p(\mu|_{A_j})\) 中的函数延拓为 \(A_j\) 外的零函数，就得到 \(L^p(\mu)\) 的子空间；\(\Lambda\) 在其上的限制范数不超过 \(\|\Lambda\|\)。前一引理给出密度 \(g_j\in L^1(\mu|_{A_j})\)。在每个 \(A_j\) 上令 \(g=g_j\)，即可得到可测函数 \(g\)，并对每个有界、在某个 \(E_N\) 外为零的 \(v\) 有
\[
 \Lambda(v)=\int_Xvg\,\mathrm d\mu.
\]
这是有限多个局部等式之和，尚未使用 \(g\in L^q(\mu)\)。

选各局部密度的几乎处处有限代表，并在相应零集上置零。令
\[
 B_N=E_N\cap\{\,|g|\le N\,\},
 \quad v_N=\vartheta_g|g|^{q-1}\mathbf1_{B_N}.
\]
函数 \(v_N\) 有界且支集测度有限，所以刚才的等式适用。记
\(a_N=\int_{B_N}|g|^q\,\mathrm d\mu<\infty\)，则
\[
 a_N=\Lambda(v_N)
 \le\|\Lambda\|\|v_N\|_p
 =\|\Lambda\|a_N^{1/p}.
\]
若 \(a_N>0\)，相除得到 \(a_N^{1/q}\le\|\Lambda\|\)；若 \(a_N=0\)，这个结论也成立。因为 \(B_N\uparrow X\)，单调收敛定理给出
\[
 \int_X|g|^q\,\mathrm d\mu\le\|\Lambda\|^q.
\]
故 \(g\in L^q(\mu)\)。

最后，对任意 \(f\in L^p(\mu)\)，取
\(f_N=f\mathbf1_{E_N\cap\{|f|\le N\}}\)。控制收敛给出 \(\|f_N-f\|_p\rightarrow0\)，而 Hölder 不等式给出
\[
 \left|\int_X(f_N-f)g\,\mathrm d\mu\right|
 \le\|f_N-f\|_p\|g\|_q\longrightarrow0.
\]
在 \(\Lambda(f_N)=\int f_Ng\,\mathrm d\mu\) 中取极限，就得到 \(\Lambda=\Lambda_g\)。\cref{b1:prop:lp-dual-isometry}给出唯一性与等范数。
\end{proof}

证明中的截断有两个不同用途：有限测度支集让指示函数进入 \(L^p\)，幅度截断则让密度所决定的相位测试在尚未知道 \(g\in L^q\) 时也能使用。若直接代入未截断的 \(\vartheta_g|g|^{q-1}\)，就会预先假定正在证明的可积性。

\subsection{\texorpdfstring{\(L^1\)}{L1} 的对偶}
\label{b1:subsec:100303}

当 \(p=1\) 时，共轭指数为无穷。局部密度仍由 Radon--Nikodym 定理产生，但有界性不再表现为某个幂的可积性，而是密度的本质上界。

\begin{theorem}[\(L^1\) 对偶表示]\label{b1:thm:l1-dual-sigma-finite}
设 \(\mu\) 为 \(\sigma\)-有限正测度，则
\[
 L^\infty(\mu;\mathbb K)\longrightarrow(L^1(\mu;\mathbb K))^*,
 \quad g\longmapsto\Lambda_g,\quad
 \Lambda_g(f)=\int_Xfg\,\mathrm d\mu
\]
是线性等距满射。
\end{theorem}
\begin{proof}
给定 \(\Lambda\in(L^1)^*\)，用\cref{b1:lem:lp-functional-local-density}在上一证明的有限测度块 \(A_j\) 上构造并拼接 \(g\)。所得 \(g\) 在每个 \(E_N\) 上可积，且表示等式对所有有界、在 \(E_N\) 外为零的测试成立。

固定 \(t>\|\Lambda\|\)，令 \(B=E_N\cap\{|g|>t\}\)。若 \(\mu(B)>0\)，取 \(v=\vartheta_g\mathbf1_B\)，便有
\[
 t\mu(B)\le\int_B|g|\,\mathrm d\mu
 =\Lambda(v)\le\|\Lambda\|\mu(B),
\]
矛盾。因此 \(\mu(E_N\cap\{|g|>t\})=0\) 对所有 \(N\) 成立。先合并这些零集，再取 \(t=\|\Lambda\|+1/k\)，可得 \(|g|\le\|\Lambda\|\) 几乎处处。故 \(g\in L^\infty\)。使用 \(L^1\) 截断逼近及
\[
 \left|\int_X(f_N-f)g\,\mathrm d\mu\right|
 \le\|f_N-f\|_1\|g\|_\infty
\]
即可把表示推广到所有 \(f\in L^1\)。

对给定 \(g\in L^\infty\)，积分估计给出 \(\|\Lambda_g\|\le\|g\|_\infty\)。若 \(0<t<\|g\|_\infty\)，集合 \(\{|g|>t\}\) 测度为正；由 \(\sigma\)-有限性，其中有一个可测子集 \(A\) 满足 \(0<\mu(A)<\infty\)。取 \(f=\vartheta_g\mathbf1_A/\mu(A)\)，便有 \(\|f\|_1=1\) 和 \(\Lambda_g(f)\ge t\)。令 \(t\uparrow\|g\|_\infty\)，得到等范数；零范数情形直接成立。等距性给出唯一性。
\end{proof}

这里的有限测度测试不能无条件找到。例如在单点空间 \(\{a\}\) 上令 \(\mu(\{a\})=\infty\)，则 \(L^1(\mu)=\{0\}\)，而 \(L^\infty(\mu)=\mathbb K\)。所有密度都在零空间上诱导零泛函，积分映射并非等距单射。这说明删除测度假设后，不能照抄端点的表示结论。

\subsection{\texorpdfstring{\(L^\infty\)}{L∞} 的特殊性}
\label{b1:subsec:100304}

对任意正测度，\(g\in L^1(\mu)\) 仍产生 \(L^\infty(\mu)\) 上的泛函，而且
\[
 \|\Lambda_g\|=\|g\|_1.
\]
上界由积分估计得到，下界由测试函数 \(\vartheta_g\in L^\infty\)、\(\|\vartheta_g\|_\infty\le1\) 得到。然而这个等距嵌入通常不满射，即使测度已经 \(\sigma\)-有限。

在 \(\mathbb N\) 上取计数测度，则 \(L^\infty=\ell^\infty\)、\(L^1=\ell^1\)。令 \(c\) 为收敛序列组成的子空间，其上的极限泛函 \(x\mapsto\lim_nx_n\) 连续且范数为 \(1\)。由 Hahn--Banach 定理取得保范延拓 \(L\in(\ell^\infty)^*\)；实情形的进一步性质见\cref{b1:ex:08-positive-limit}，复情形使用复线性保范延拓即可。

这个 \(L\) 不能由某个 \(a\in\ell^1\) 表示。否则假设
\[
 L(x)=\sum_{n=1}^{\infty}a_nx_n.
\]
代入每个单位向量 \(e_n\)，由其通常极限为零可得 \(a_n=L(e_n)=0\)；但代入常数序列 \(\mathbf1\)，又有 \(L(\mathbf1)=1\)，矛盾。

这一例子也指出 Radon--Nikodym 方法在这里的断点。令 \(\nu(E)=L(\mathbf1_E)\)，则 \(\nu\) 有限可加，却有
\[
 \nu(\mathbb N)=1,\quad \nu(\{n\})=0\quad(n\ge1),
\]
所以不可数可加。部分和 \(\sum_{n=1}^N e_n\) 并不在 \(\ell^\infty\) 范数下趋于 \(\mathbf1\)，泛函连续性因而不能提供前面构造测度时的极限等式。这里使用的只是通常极限泛函的延拓，没有额外要求移位不变性。

% !TeX root = ../../Book1.tex
\section{反身性与弱收敛}
\label{b1:sec:1004}

对偶表示把 \(L^p\) 上抽象的连续泛函变成了乘以函数再积分的操作。将这个表示使用两次，就能判断自然嵌入是否满射；第\ref{b1:ch:09}章的弱紧性理论也随之适用。以下除明确放宽条件的命题外，均设 \((X,\mathcal A,\mu)\) 为 \(\sigma\)-有限正测度空间，并固定 \(1<p<\infty\)、\(q=p'\)。

% 来源：Axler2020Measure，7B习题19及9B习题17提出双对偶反身性；
% 本节用已证sigma有限对偶表示完整展开自然嵌入计算。
% 弱子列与凸泛函方法调用本书第9章，不把一般测度对偶版本当作已证。

\subsection{\texorpdfstring{\(1<p<\infty\)}{1<p<∞} 的反身性}
\label{b1:subsec:100401}

“对偶的对偶又是原来的函数空间”还不足以单独证明反身性；必须核对两次表示的复合确实是自然嵌入。上一节采用的双线性积分配对使这个核对十分直接。

\begin{theorem}\label{b1:thm:lp-reflexive}
设 \(\mu\) 为 \(\sigma\)-有限正测度，\(1<p<\infty\)。则实或复空间 \(L^p(\mu)\) 反身。
\end{theorem}
\begin{proof}
由\cref{b1:thm:lp-dual}，映射 \(D_p:g\mapsto\Lambda_g\) 是从 \(L^q(\mu)\) 到 \(\bigl(L^p(\mu)\bigr)^*\) 的线性等距满射。任取
\(\Phi\in\bigl(L^p(\mu)\bigr)^{**}\)，令
\[
 \Psi(g)=\Phi(D_pg)\quad\bigl(g\in L^q(\mu)\bigr).
\]
因为 \(|\Psi(g)|\le\|\Phi\|\cdot\|g\|_q\)，所以 \(\Psi\in\bigl(L^q(\mu)\bigr)^*\)。指数 \(q\) 也在 \((1,\infty)\) 中，其共轭指数为 \(p\)。再次应用对偶表示，取得 \(f\in L^p(\mu)\)，使
\[
 \Phi(D_pg)=\Psi(g)=\int_Xgf\,\mathrm d\mu
 =(D_pg)(f)\quad\bigl(g\in L^q(\mu)\bigr).
\]
由于每个 \(L^p\) 上的连续泛函都等于某个 \(D_pg\)，上式就是
\(\Phi(\Lambda)=\Lambda(f)\) 对所有 \(\Lambda\in(L^p)^*\) 成立。因此 \(\Phi=J_{L^p}f\)，自然嵌入满射。
\end{proof}

指数范围在第二次使用对偶表示时至关重要。若从 \(p=1\) 出发，第一次得到 \(L^\infty\)，第二次便遇到上一节的额外泛函，不能再写回 \(L^1\)。具体地，\((\ell^1)^*=\ell^\infty\)，而极限延拓泛函 \(L\in(\ell^\infty)^*\) 不由任何 \(\ell^1\) 元素表示，故不在 \(J_{\ell^1}(\ell^1)\) 中。因此 \(\ell^1\) 不反身。

同样，\(\ell^\infty\) 含有闭线性子空间 \(c_0\)，而第\ref{b1:sec:0904}节已说明 \(c_0\) 不反身；由\cref{b1:cor:closed-subspace-reflexive}，\(\ell^\infty\) 也不反身。这不意味着每个 \(L^1\) 或 \(L^\infty\) 空间都不反身：有限集合上各点质量为正且有限时，它们都有限维。端点所缺少的是对一般测度空间统一成立的反身性结论。

\subsection{有界序列的弱收敛子列}
\label{b1:subsec:100402}

\begin{corollary}\label{b1:cor:lp-weak-subsequence}
设 \(\sup_n\|f_n\|_p<\infty\)。则存在子列 \((f_{n_k})\) 及 \(f\in L^p(\mu)\)，使
\[
 \int_X f_{n_k}g\,\mathrm d\mu
 \longrightarrow\int_Xfg\,\mathrm d\mu
 \quad\bigl(g\in L^q(\mu)\bigr).
\]
这等价于 \(f_{n_k}\rightharpoonup f\) 在 \(L^p(\mu)\) 中成立。
\end{corollary}
\begin{proof}
\cref{b1:thm:lp-reflexive}与\ref{b1:thm:reflexive-weak-subsequence}给出弱收敛子列；\cref{b1:thm:lp-dual}说明全部连续线性泛函恰是所列积分，所以弱收敛等价于这些积分收敛。
\end{proof}

此结论不要求 \(L^p(\mu)\) 可分，也没有断言函数值逐点收敛。它适合已知范数估计、却尚未找到点态控制的近似过程。由范数的弱下半连续性，还能保留估计
\[
 \|f\|_p\le\liminf_{k\to\infty}\|f_{n_k}\|_p.
\]

\begin{proposition}\label{b1:prop:lp-weak-finite-set-tests}
设 \(f_n,f\in L^p(\mu)\)，且 \(\sup_n\|f_n\|_p<\infty\)。则 \(f_n\rightharpoonup f\)，当且仅当对每个有限测度可测集 \(E\)，有
\[
 \int_E f_n\,\mathrm d\mu\longrightarrow\int_E f\,\mathrm d\mu.
\]
\end{proposition}
\begin{proof}
若弱收敛成立，取测试函数 \(\mathbf1_E\in L^q(\mu)\) 即得结论。

反之，给定条件及有限线性组合说明积分收敛对所有有限测度支集简单函数成立。由\cref{b1:prop:lp-simple-dense}，这些函数在 \(L^q(\mu)\) 中稠密。记 \(C=\sup_n\|f_n\|_p+\|f\|_p\)。对 \(g\in L^q\) 和这样的简单函数 \(s\)，Hölder 不等式给出
\[
 \left|\int_X(f_n-f)g\,\mathrm d\mu\right|
 \le\left|\int_X(f_n-f)s\,\mathrm d\mu\right|
   +C\|g-s\|_q.
\]
先选择 \(s\) 使第二项足够小，再令 \(n\) 足够大，即得所有 \(g\) 上的收敛。
\end{proof}

这个判别仍需要统一范数界。它把测试函数从整个 \(L^q\) 简化为集合的指示函数，并没有仅凭集合积分的逐点收敛制造范数估计。

作为对照，在 \(\ell^p\) 中，单位向量 \(e_n\) 对每个 \(g\in\ell^q\) 满足
\(\Lambda_g(e_n)=g_n\rightarrow0\)，因为 \(q<\infty\) 保证 \(\ell^q\) 序列的各项趋零。所以 \(e_n\rightharpoonup0\)，尽管 \(\|e_n\|_p=1\)，任意两个不同单位向量的距离都是 \(2^{1/p}\)。弱紧性因此不会恢复范数紧性。

若换成 \(\ell^1\)，所有坐标评价仍使 \(e_n\) 趋向零，但对偶中的常数序列 \(\mathbf1\in\ell^\infty\) 给出 \(\Lambda_{\mathbf1}(e_n)=1\)。任何弱收敛子列的极限又都会被坐标评价强制为零，因此这里根本没有弱收敛子列。对偶指数从有限变为无穷，使有限支撑测试不再足以逼近所有测试函数。

\subsection{凸泛函与弱下半连续性}
\label{b1:subsec:100403}

对积分能量，先在范数收敛下使用逐点工具，再借凸性转到弱拓扑，往往比直接处理弱收敛函数值更方便。

\begin{proposition}[凸积分泛函的弱下半连续性]\label{b1:prop:convex-integral-weak-lsc}
设 \(\mu\) 为任意正测度，\(1\le r<\infty\)。设
\(\phi:\mathbb K\rightarrow[0,\infty]\) 下半连续、凸，且 \(\phi(0)=0\)；复情形把 \(\mathbb C\) 看作实平面理解凸性。则
\[
 \mathcal F(f)=\int_X\phi\bigl(f(x)\bigr)\,\mathrm d\mu(x)
 \quad\bigl(f\in L^r(\mu;\mathbb K)\bigr)
\]
是强下半连续、弱下半连续的凸泛函。
\end{proposition}
\begin{proof}
下半连续函数 \(\phi\) 可测，所以非负积分有定义，且不随 \(f\) 的几乎处处代表改变。\(\phi(0)=0\) 保证 \(\mathcal F(0)=0\)，故泛函不恒等于无穷。逐点凸不等式经积分得到 \(\mathcal F\) 的凸性。

设 \(f_n\rightarrow f\) 在 \(L^r\) 范数下成立。若 \(\liminf_n\mathcal F(f_n)=\infty\)，所需不等式直接成立。否则先选子列，使其能量趋于这个有限下极限，再取更稀的子列 \((f_{n_k})\)，使
\[
 \sum_{k=1}^{\infty}\|f_{n_k}-f\|_r^r<\infty.
\]
由非负积分的单调收敛定理，
\[
 \int_X\sum_{k=1}^{\infty}|f_{n_k}-f|^r\,\mathrm d\mu
 =\sum_{k=1}^{\infty}\|f_{n_k}-f\|_r^r<\infty.
\]
因此逐点级数几乎处处有限，特别地 \(f_{n_k}\rightarrow f\) 几乎处处。下半连续性与 Fatou 引理给出
\[
 \mathcal F(f)\le
 \int_X\liminf_k\phi(f_{n_k})\,\mathrm d\mu
 \le\liminf_k\mathcal F(f_{n_k})
 =\liminf_n\mathcal F(f_n).
\]
这证明了范数拓扑下的下半连续性。再由\cref{b1:thm:convex-weak-lsc}，它也是弱下半连续的。
\end{proof}

例如取 \(\phi(z)=|z|^s\)、\(1\le s<\infty\)，可知即使在所讨论的 \(L^r\) 上积分可能为无穷，这个 \(s\) 次能量仍弱下半连续。这里没有把弱收敛误作几乎处处收敛：逐点子列只用于证明强下半连续，转到弱拓扑时另用了凸性。

回到本节的 \(\sigma\)-有限测度与 \(1<p<\infty\)，取 \(g_1,\ldots,g_m\in L^q\) 及标量 \(a_1,\ldots,a_m\)，假设约束集
\[
 \mathcal C=\left\{\,f\in L^p\ \middle|\
   \int_Xfg_j\,\mathrm d\mu=a_j,\ 1\le j\le m\,\right\}
\]
非空。每个约束都是弱连续泛函的水平集，所以 \(\mathcal C\) 弱闭。选取 \(\|f\|_p^p\) 在 \(\mathcal C\) 上的极小化序列，其范数有界；由\cref{b1:cor:lp-weak-subsequence}取得弱收敛子列。极限仍满足全部约束，而弱下半连续性保证它达到下确界。这是\cref{b1:thm:direct-method}在积分约束下的具体实现。

单个非零约束函数 \(g\in L^q\) 时，答案还可以直接计算。若要求 \(\int fg\,\mathrm d\mu=a\)，Hölder 不等式给出 \(\|f\|_p\ge |a|/\|g\|_q\)；而
\[
 f_0=\dfrac{a\,\vartheta_g|g|^{q-1}}{\|g\|_q^q}
\]
满足约束，并有 \(\|f_0\|_p=|a|/\|g\|_q\)。抽象的存在性过程与对偶范数中达到等号的相位测试，在这个例子中给出同一个极小点。

% !TeX root = ../../Book1.tex
\OptionalSection{\texorpdfstring{\(L^1\)}{L¹} 的特殊性}{b1:sec:1005}

在 \(1<p<\infty\) 时，有界性足以保证从 \(L^p\) 序列中抽取弱收敛子列。端点 \(p=1\) 的困难在于集中：积分总量可以保持有界，甚至恒定，而越来越多的质量落入越来越小的集合。第五章的一致可积性正是对这种现象的限制。本节说明它如何补足 \(L^1\) 的弱紧性，并在只有有界性时保留一个较弱的子列结论。

除专门讨论无穷测度的段落外，本节固定有限测度空间 \((\Omega,\mathcal A,\mu)\)，即 \(\mu(\Omega)<\infty\)。函数可取实值或复值；涉及点值时选取有限值可测代表。

% 来源：Fremlin2016Measure，卷2，246A、246G(iii)、247C及247E，
% 已读取作者官网 chap24.pdf 的相应正文。充分性用双对偶弱-*闭包和RN重述；
% 必要性明确保留proofsketch，不把互不相交化与滑动选择装成完整证明。
% BallMurat1989Biting：已核作者出版清单、正式元数据与作者公开摘要，
% 未读取原文证明；下文标量证明自行展开幅值尾部的对角构造。

\subsection{一致可积性再论}
\label{b1:subsec:100501}

沿用\cref{b1:def:uniform-integrability}，非空族 \(\mathcal F\subset L^1(\mu)\) 一致可积是指
\[
  \lim_{M\to\infty}\sup_{f\in\mathcal F}
  \int_{\{|f|>M\}}|f|\,\mathrm d\mu=0.
\]
由\cref{b1:prop:uniform-integrability-criterion}，在当前有限测度假设下，这等价于 \(L^1\) 范数一致有界与积分的一致绝对连续性同时成立。为便于后文使用，记
\[
  B=\sup_{f\in\mathcal F}\|f\|_1,
  \quad
  \omega(\delta)=\sup_{\substack{f\in\mathcal F,\ A\in\mathcal A\\
                               \mu(A)\leq\delta}}
                 \int_A|f|\,\mathrm d\mu.
\]
于是上述等价条件就是 \(B<\infty\) 且 \(\omega(\delta)\to0\) 当 \(\delta\downarrow0\)。这里的 \(\omega\) 控制所有小集合，而不只控制某个预先选定的集中位置。

一个稍强的可积性估计往往能给出这种控制。若 \(1<p<\infty\)，且 \(\sup_{f\in\mathcal F}\|f\|_p\leq C\)，则
\[
  \int_{\{|f|>M\}}|f|\,\mathrm d\mu
  \leq M^{1-p}\int_\Omega|f|^p\,\mathrm d\mu
  \leq C^pM^{1-p}\longrightarrow0
\]
一致地对 \(f\) 成立。有限测度又保证这些函数属于 \(L^1\)。不过，一致可积性比“存在某个统一的 \(p>1\) 上界”更灵活：它只要求幅值尾部能够统一趋零，不指定幂次速度。

有限测度这一限定不能在引用定理时略去。在一般 \(\sigma\)-有限空间上，本书的幅值尾部条件还须补上空间尾部条件
\begin{equation}
\label{b1:eq:l1-no-escape}
  \text{对每个 }\varepsilon>0\text{，存在 }E\in\mathcal A,
  \quad\mu(E)<\infty,\quad
  \sup_{f\in\mathcal F}\int_{\Omega\setminus E}|f|\,\mathrm d\mu
  <\varepsilon.
\end{equation}
例如，在 \(\mathbb R\) 上令 \(f_n=\mathbf1_{[n,n+1]}\)。它们有统一幅值上界和 \(L^1\) 范数，却不满足\eqref{b1:eq:l1-no-escape}：对任意有限测度集 \(E\)，由区间互不相交可知 \(\mu(E\cap[n,n+1])\to0\)。所以集外积分趋于一，质量没有集中在小集合上，而是向无穷远移动。

本书的幅值尾部条件与\eqref{b1:eq:l1-no-escape}合起来，等价于 Fremlin 在一般测度空间上的一致可积性定义\cite[第2卷，246A]{Fremlin2016Measure}。一个方向由
\[
  \int_\Omega\bigl(|f|-M\mathbf1_E\bigr)_+\,\mathrm d\mu
  \leq\int_{\{|f|>M\}}|f|\,\mathrm d\mu
       +\int_{\Omega\setminus E}|f|\,\mathrm d\mu
\]
给出；反向用 \(|f|\leq2(|f|-M)_+\) 在 \(\{|f|>2M\}\) 上成立，以及集外部分原封不动地出现在左侧积分中。两项控制也给出 \(L^1\) 有界性，因为剩余部分的积分不超过 \(M\mu(E)\)。以下主定理先只在有限测度空间陈述，避免混用定义。

\subsection{Dunford–Pettis 定理}
\label{b1:subsec:100502}

回顾\cref{b1:def:weakly-compact-set}，相对弱紧要求弱闭包是紧集，并不要求原集合本身闭。\term*[dunford-pettis]{Dunford--Pettis 定理}[Dunford--Pettis theorem]把这个拓扑条件化为积分条件。

\begin{theorem}[Dunford--Pettis，有限测度情形]
\label{b1:thm:dunford-pettis}
设 \(\mu(\Omega)<\infty\)，\(\mathcal F\subset L^1(\mu)\) 非空。则 \(\mathcal F\) 相对弱紧，当且仅当 \(\mathcal F\) 一致可积；等价地，当且仅当它 \(L^1\) 有界且积分一致绝对连续。
\end{theorem}

下面完整证明从一致可积性到相对弱紧性。反向所需的集合选择较长，将另列证明概要。一般测度空间的正式表述及完整证明见 Fremlin 的《Measure Theory》\cite[第2卷，247C；复值情形见247E]{Fremlin2016Measure}；使用其一般版本时应采用上一小节说明的两项尾部条件。

\begin{proof}[充分性的证明]
设 \(\mathcal F\) 一致可积，取前述 \(B\) 与 \(\omega\)。每个 \(f\in L^1\) 定义 \(L^\infty\) 上的连续线性泛函
\[
  T_f(h)=\int_\Omega fh\,\mathrm d\mu,
  \quad h\in L^\infty.
\]
有 \(\|T_f\|\leq\|f\|_1\)。反向在 \(f\ne0\) 处取 \(h=\overline f/|f|\)，在其余点取零，即得 \(\|T_f\|=\|f\|_1\)。由\cref{b1:thm:banach-alaoglu}，\(\{T_f:f\in\mathcal F\}\) 在 \((L^\infty)^*\) 中的弱-*闭包 \(K\) 是紧集，且其中每个泛函的范数不超过 \(B\)。

关键是证明 \(K\) 不会新增无法由 \(L^1\) 函数表示的泛函。固定 \(T\in K\)，令
\[
  \nu(A)=T(\mathbf1_A),\quad A\in\mathcal A.
\]
由线性性，\(\nu\) 有限可加。对固定的 \(A\)，泛函的取值映射在弱-*拓扑下连续，所以
\begin{equation}
\label{b1:eq:dp-limit-measure-control}
  |\nu(A)|\leq\sup_{f\in\mathcal F}
                  \int_A|f|\,\mathrm d\mu
              \leq\omega(\mu(A)).
\end{equation}
若 \(A_n\downarrow\varnothing\)，有限测度保证 \(\mu(A_n)\to0\)，从而 \(\nu(A_n)\to0\)。有限可加性与这个从上连续性合起来给出可数可加性：对互不相交的 \(A_j\)，把并集减去前 \(n\) 项所余的集合递减到空集即可。

此外，\(\nu\) 的全变差有限。对任意有限可测分割 \((A_j)\)，选取 \(|c_j|\leq1\) 使 \(c_j\nu(A_j)=|\nu(A_j)|\)，便有
\[
  \sum_j|\nu(A_j)|
  =T\left(\sum_jc_j\mathbf1_{A_j}\right)\leq B.
\]
因此 \(|\nu|(\Omega)\leq B\)。由\eqref{b1:eq:dp-limit-measure-control}还知 \(\nu\ll\mu\)。实值情形可用\cref{b1:cor:signed-radon-nikodym}；复值情形分别对 \(\operatorname{Re}\nu\) 与 \(\operatorname{Im}\nu\) 使用该结论。于是存在 \(g\in L^1(\mu)\)，使 \(\nu(A)=\int_Ag\,\mathrm d\mu\)。

从示性函数到简单函数用线性性，再用简单函数在 \(L^\infty\) 范数下的逼近，得到
\[
  T(h)=\int_\Omega gh\,\mathrm d\mu\quad(h\in L^\infty).
\]
故 \(T=T_g\)，即 \(K\) 中每个点仍来自 \(L^1\)。最后由\cref{b1:thm:l1-dual-sigma-finite}，映射 \(f\mapsto T_f\) 把 \(L^1\) 的弱拓扑恰好变为像上的弱-*拓扑，因为两边都由全部积分 \(\int fh\,\mathrm d\mu\)、\(h\in L^\infty\) 决定。紧集 \(K\) 因而对应于 \(L^1\) 中包含 \(\mathcal F\) 的弱紧集。这证明了相对弱紧性。
\end{proof}

\begin{proofsketch}
必要性先看实值情形。相对弱紧族在每个连续线性泛函下的像有界，由一致有界性原理可得 \(L^1\) 范数有界。余下的实质是排除积分在互不相交的集合上持续保留同等大小的质量。

所需的集合判据为：一个 \(L^1\) 有界族一致可积，当且仅当对每列互不相交的可测集 \((A_n)\)，都有
\[
  \sup_{f\in\mathcal F}\left|\int_{A_n}f\,\mathrm d\mu\right|
  \longrightarrow0.
\]
这里绝对值放在积分外；把这个条件转回 \(\int_A|f|\) 的一致小量，需要实质性的递归选集论证，见\cite[第2卷，246G(iii)]{Fremlin2016Measure}。

若上述条件失败，便可取 \(f_n\in\mathcal F\) 与互不相交的 \(A_n\)，使相应积分的绝对值均不小于某个 \(\varepsilon>0\)。弱紧性提供聚点。依次抽取函数和集合，可以使每个新函数在先前选定集合上的积分接近聚点的积分，同时使先前函数在以后选定集合上的积分总量很小。将选出的集合取并，就得到一个固定的有界测试函数；它在所选函数上的积分与任何候选聚点上的积分相冲突，从而否定弱紧性。这一选择及误差求和的完整步骤见\cite[第2卷，247C(a)]{Fremlin2016Measure}。本概要不代替上述两处的递归证明。

复值情形将族映到实部族与虚部族；这些映射对于相应的弱拓扑连续，故两个像仍相对弱紧。对它们应用实值结论，再用 \(|f|\leq|\operatorname{Re}f|+|\operatorname{Im}f|\) 和一致绝对连续性判据，即得原族一致可积。
\end{proofsketch}

\subsection{\texorpdfstring{\(L^1\)}{L1} 弱紧性}
\label{b1:subsec:100503}

\begin{corollary}[一致可积列的弱收敛子列]
\label{b1:cor:l1-ui-weak-subsequence}
若 \(\mu(\Omega)<\infty\)，且 \((f_n)\subset L^1(\mu)\) 一致可积，则存在子列 \((f_{n_j})\) 与 \(f\in L^1(\mu)\)，使
\[
  \int_\Omega f_{n_j}h\,\mathrm d\mu
  \longrightarrow\int_\Omega fh\,\mathrm d\mu
  \quad(h\in L^\infty(\mu)).
\]
\end{corollary}
\begin{proof}
由刚证的充分性，\(\{f_n:n\geq1\}\) 的弱闭包紧。应用\cref{b1:thm:eberlein-smulian}中弱紧性蕴含子列性质的方向，再用 \(L^1\) 对偶表示即可。这里使用的是第九章已经展开证明的方向，不需要该处仅列概要的逆向。
\end{proof}

因此，若 \(C\subset L^1\) 非空、凸、范数闭且一致可积，则 \(C\) 弱紧：闭凸性给出弱闭性，Dunford--Pettis 定理给出相对弱紧性。这使第九章的直接法仍有用武之地。对于 \(C\) 上有下界且至少在一点有限的弱下半连续泛函，选取极小化序列，抽取弱收敛子列，再比较极限处的泛函值，就能得到极小点。这里不需要假设整个 \(L^1\) 空间反身；提供弱紧性的是具体可行族的一致可积性。

第五章的集中尖峰 \(f_n=n\mathbf1_{(0,1/n)}\) 没有弱收敛子列。事实上，若某个子列弱收敛于 \(f\)，则对任意 \(\delta>0\) 及可测集 \(A\subset(\delta,1)\)，测试 \(\mathbf1_A\) 得到 \(\int_Af=0\)，故 \(f=0\) 几乎处处于 \((\delta,1)\)。让 \(\delta\downarrow0\) 得 \(f=0\)，但测试常函数一又给出 \(\int_0^1f=1\)，矛盾。这也说明，不能把单位球有界当作 \(L^1\) 中的弱紧性证明。

另一方面，弱紧性仍不等于范数紧性。在 \((0,1)\) 上令
\[
  r_n(x)=(-1)^{\lfloor2^nx\rfloor},\quad n\geq1.
\]
这些函数均只有 \(\pm1\) 两个值，故一致可积。它们在 \(L^2(0,1)\) 中构成标准正交系：当 \(m>n\) 时，在 \(r_n\) 为常数的每个二进分割区间内，\(r_m\) 取正负值的长度相等。因此，对每个 \(h\in L^\infty(0,1)\subset L^2(0,1)\)，Bessel 不等式给出 \(\int r_nh\to0\)，即 \(r_n\rightharpoonup0\) 于 \(L^1\)。然而对 \(m\ne n\)，
\[
  \|r_m-r_n\|_1=1,
\]
所以不存在范数收敛子列。一致可积性排除了质量集中，并没有排除这样的振荡。

\subsection{Chacon 引理}
\label{b1:subsec:100504}

若只能得到 \(L^1\) 有界性，未必能在整个空间上取得弱极限。不过可以先去掉测度任意小的坏集。\term[chacon-yinli]{Chacon 引理}[Chacon's biting lemma]说明，经过抽取子列，这些坏集可以选成一个递减的族，并使每个固定坏集之外都恢复一致可积性。有关标准表述及向量值推广，可参见 Ball 与 Murat 的论文\cite{BallMurat1989Biting}。下面给出实值与复值函数均适用的尾部构造。

\begin{theorem}[Chacon 引理，有限测度情形]
\label{b1:thm:chacon-biting}
设 \(\mu(\Omega)<\infty\)，\((f_n)\subset L^1(\mu)\)，且 \(\sup_n\|f_n\|_1\leq B<\infty\)。则存在子列 \((f_{n_j})\) 和递减的可测集 \(E_k\)，使
\[
  \mu(E_k)\longrightarrow0,
  \quad
  \{f_{n_j}\mathbf1_{\Omega\setminus E_k}:j\geq1\}
  \text{ 对每个固定的 }k\text{ 一致可积}.
\]
还可以进一步抽取子列并找到 \(f\in L^1(\mu)\)，使对每个固定的 \(k\)，
\[
  f_{n_j}\mathbf1_{\Omega\setminus E_k}
  \rightharpoonup f\mathbf1_{\Omega\setminus E_k}
  \quad\text{于 }L^1(\mu).
\]
\end{theorem}

\begin{proof}
对正整数 \(m\) 记
\[
  a_n(m)=\int_{\{|f_n|>m\}}|f_n|\,\mathrm d\mu.
\]
每个数列 \((a_n(m))_n\) 都位于 \([0,B]\)。逐个 \(m\) 抽取子列，再取对角子列，可以假定对所有正整数 \(m\)，都有 \(a_n(m)\to a(m)\)。由于 \(a(m)\) 随 \(m\) 递减且非负，存在
\[
  \alpha=\lim_{m\to\infty}a(m)\in[0,B].
\]

选取严格递增的正整数 \(M_j\geq2^j(B+1)\)，再选取严格递增的指标 \(n_j\)，使
\[
  |a_{n_j}(m)-a(m)|<\frac1j
  \quad(1\leq m\leq M_j).
\]
每一步只有有限多个收敛条件，故这样的选择总能完成。于是 \(a_{n_j}(M_j)\to\alpha\)。令
\[
  A_j=\{|f_{n_j}|>M_j\},\quad
  g_j=f_{n_j}\mathbf1_{\Omega\setminus A_j}.
\]
我们先证明 \((g_j)\) 一致可积。固定正整数 \(m\)，当 \(j\) 充分大而 \(M_j\geq m\) 时，有
\[
  \int_{\{|g_j|>m\}}|g_j|\,\mathrm d\mu
  =a_{n_j}(m)-a_{n_j}(M_j)
  \longrightarrow a(m)-\alpha.
\]
给定 \(\varepsilon>0\)，先取 \(m\) 使 \(a(m)-\alpha<\varepsilon/2\)，上述收敛便保证除有限多个 \(j\) 外，尾部积分都小于 \(\varepsilon\)。对余下的有限多个可积函数再增大截断高度，使它们的尾部积分也小于 \(\varepsilon\)。增大高度不会破坏已经得到的估计，所以全族一致可积。

现在令
\[
  E_k=\bigcup_{j\geq k}A_j.
\]
由 Markov 不等式，
\[
  \mu(E_k)\leq\sum_{j\geq k}\frac{B}{M_j}
  \leq\sum_{j\geq k}2^{-j}=2^{1-k}\longrightarrow0.
\]
对固定的 \(k\) 与所有 \(j\geq k\)，\(A_j\subset E_k\)，因而
\[
  f_{n_j}\mathbf1_{\Omega\setminus E_k}
  =g_j\mathbf1_{\Omega\setminus E_k}.
\]
限制到一个可测集不会增加幅值尾部积分，故右侧构成一致可积族。再加上 \(j<k\) 的有限多个函数，仍一致可积，第一部分得证。

最后证明可取相容的弱极限。记 \(F_k=\Omega\setminus E_k\)，则 \(F_k\) 递增，且其并覆盖 \(\Omega\) 除去一个零集。对 \(k=1,2,\ldots\) 依次应用\cref{b1:cor:l1-ui-weak-subsequence}，再抽取对角子列，便可假定对每个固定的 \(k\)，\(f_{n_j}\mathbf1_{F_k}\) 弱收敛于某个 \(h_k\in L^1\)。

乘以 \(\mathbf1_{F_k}\) 是弱连续的有界线性算子，因为每个测试函数 \(h\in L^\infty\) 仍被变为 \(h\mathbf1_{F_k}\in L^\infty\)。所以 \(h_k\) 支持在 \(F_k\) 上，且当 \(k<\ell\) 时，弱极限的唯一性给出 \(h_\ell\mathbf1_{F_k}=h_k\)。在去掉可数个例外零集后，可以把它们拼成一个可测函数 \(f\)，使 \(f\mathbf1_{F_k}=h_k\)。范数的弱下半连续性给出
\[
  \int_{F_k}|f|\,\mathrm d\mu=\|h_k\|_1
  \leq\liminf_{j\to\infty}\|f_{n_j}\mathbf1_{F_k}\|_1\leq B.
\]
令 \(k\to\infty\)，单调收敛定理便给出 \(\|f\|_1\leq B\)。这证明了所需的弱收敛结论。
\end{proof}

这里固定坏集和改变坏集是两回事。结论说的是：先固定 \(k\)，再令子列指标 \(j\to\infty\)。它并不保证 \(\int_{E_k}|f_{n_j}|\) 随 \(k\) 一致趋零，也就不能直接把坏集补回去。对集中尖峰 \(f_n=n\mathbf1_{(0,1/n)}\)，取 \(E_k=(0,1/k)\)，则每个 \(E_k\) 外的序列最终恒为零；但全空间积分始终等于一。Chacon 引理恢复了坏集外的弱极限，可能丢失的集中质量仍须由问题本身提供额外估计。

% 本章小结（不计入编号）
\begin{chaptersummary}
\(L^p\) 的元素是几乎处处等价类。Hölder 不等式控制乘积积分，
Minkowski 不等式保证范数的三角关系，可求和的误差列则把积分估计转成空间的完备性。
有限 \(p\) 时可以丢掉很高的值以及很小的非零值，再作简单函数逼近；
\(p=\infty\) 时误差按本性上确界衡量，丢掉小测度集合不一定使误差变小。

对偶表示把连续线性泛函写成积分测试。对 \(1<p<\infty\)，
共轭指数再次取共轭后返回原指数，因此自然嵌入成为满射，自反性随之成立。
这为有界列提供弱收敛子列，并使闭凸约束上的极小化方法能够直接用于函数空间。
在本章采用的 \(\sigma\)-有限假设下，\((L^1)^*\) 可等距识别为 \(L^\infty\)，二次对偶却不再自动回到原空间。

而 \(L^1\) 有界只控制总质量，不能排除质量集中在越来越小的集合上。
有限测度空间中，一致可积性正是相对弱紧性的额外条件；
空间非有限时，还要防止质量整体移向无穷远。
Chacon 引理则允许从一般有界列中抽取子列，并移去测度趋零的递减坏集，
在其外恢复一致可积性。这些结论控制的仍是相应的弱极限或局部化极限，
不能据此直接宣称原函数列在范数下收敛。
\end{chaptersummary}


\BookExercises

% 习题以本章工具独立组织。幂函数、p趋∞和指数插值参考Axler 7A的训练主题；
% 以下解答自行推导，不逐字摘录。变分题使用本章与第9章的已证结果。

\begin{exercise}[基础运算]\label{b1:ex:10-basic-norm-tools}
设 \((X,\mathcal A,\mu)\) 为测度空间。

（1）对可测函数 \(f\) 和 \(M\geq0\)，证明
\[
 \|f\|_\infty\leq M
 \quad\Longleftrightarrow\quad
 |f|\leq M\quad\mu\text{-几乎处处}.
\]

（2）若 \(1\leq p\leq\infty\)、\(q=p'\)，证明 \(f\in L^p\)、\(g\in L^q\) 时 \(fg\in L^1\)，并写出相应范数估计；又证明 \(u,v\in L^p\) 时 \(u+v\in L^p\)。指出两步分别使用哪一个不等式。

（3）若 \(1\leq p<\infty\)、\(f\in L^p\)，对\cref{b1:subsec:100202}的幅值截断 \(T_Mf\) 证明
\[
 \|f-T_Mf\|_p^p
 =\int_{\{|f|>M\}}(|f|-M)^p\,\mathrm d\mu
 \longrightarrow0.
\]

（4）仍设 \(1\leq p<\infty\)、\(f\in L^p\)。若 \(E_n\uparrow X\)，证明 \(f\mathbf1_{E_n}\to f\) 于 \(L^p\)；再在 \(L^\infty(\mathbb R)\) 中给出同一结论失败的例子。
\end{exercise}
\begin{solution}
（1）若 \(|f|\leq M\) 几乎处处，则 \(M\) 属于本性上界的定义集合，所以 \(\|f\|_\infty\leq M\)。反之，记该向上封闭的定义集合为 \(\mathcal U\)。由 \(\inf\mathcal U\leq M<M+1/k\)，可取 \(a_k\in\mathcal U\) 使 \(a_k<M+1/k\)，再由向上封闭性得到 \(M+1/k\in\mathcal U\)。因此
\(\mu(\{|f|>M+1/k\})=0\)。取可数并得到 \(\mu(\{|f|>M\})=0\)。

（2）Hölder 不等式给出
\[
 \|fg\|_1\leq\|f\|_p\cdot\|g\|_q,
\]
所以乘积可积；端点 \((p,q)=(1,\infty)\) 与 \((\infty,1)\) 也包含在内。Minkowski 不等式给出
\[
 \|u+v\|_p\leq\|u\|_p+\|v\|_p<\infty,
\]
所以和仍属于 \(L^p\)。

（3）在 \(\{|f|\leq M\}\) 上误差为零，在其补集上误差的模等于 \(|f|-M\)，故等式成立。被积函数随 \(M\to\infty\) 逐点趋于零，并由 \(|f|^p\in L^1\) 控制；控制收敛定理给出所述极限。

（4）误差的 \(p\) 次方为 \(|f|^p\mathbf1_{X\setminus E_n}\)，它逐点趋于零且由 \(|f|^p\) 控制，故再次由控制收敛定理得到范数收敛。在 \(\mathbb R\) 上取 \(f=1\)、\(E_n=[-n,n]\)，则 \(E_n\uparrow\mathbb R\)，但
\(\|f-f\mathbf1_{E_n}\|_\infty=1\) 对所有 \(n\) 成立。
\end{solution}

\begin{exercise}\label{b1:ex:10-power-ranges}
在 Lebesgue 测度下，分别判定 \(x^{-\alpha}\mathbf1_{(0,1)}(x)\) 和
\(x^{-\beta}\mathbf1_{(1,\infty)}(x)\) 属于 \(L^p(\mathbb R)\) 的条件，其中 \(\alpha,\beta>0\)、\(1\leq p<\infty\)。
由此证明，若 \(1\leq p<q<\infty\)，则 \(L^p(\mathbb R)\) 与 \(L^q(\mathbb R)\) 互不包含。
\end{exercise}
\begin{solution}
直接计算广义积分得
\[
 \int_0^1x^{-\alpha p}\,dx<\infty\iff\alpha p<1,
 \quad \int_1^\infty x^{-\beta p}\,dx<\infty\iff\beta p>1.
\]
临界指数等于一时均按对数发散。
取 \(1/q\leq\alpha<1/p\)，第一种函数属于 \(L^p\) 而不属于 \(L^q\)；
取 \(1/q<\beta\leq1/p\)，第二种函数属于 \(L^q\) 而不属于 \(L^p\)。
局部奇异性与无穷远衰减给出的限制方向不同，不能把有限测度空间中的包含关系原样搬到整条实轴上。
\end{solution}

\begin{exercise}\label{b1:ex:10-p-to-infinity}
设 \(0<\mu(\Omega)<\infty\)，\(f\) 可测。允许范数取无穷值，证明
\(\lim_{p\to\infty}\|f\|_p=\|f\|_\infty\)。
\end{exercise}
\begin{solution}
若 \(M=\|f\|_\infty<\infty\)，则 \(\|f\|_p\leq M\mu(\Omega)^{1/p}\)，故上极限不超过 \(M\)。
对任意 \(0<c<M\)，集合 \(A_c=\{|f|>c\}\) 有正测度，且
\[
 \|f\|_p\geq c\mu(A_c)^{1/p}.
\]
取下极限得到至少为 \(c\)，再令 \(c\uparrow M\)。若 \(M=0\)，结论直接成立；
若 \(M=\infty\)，同一估计对所有 \(c>0\) 成立，故下极限为无穷。
这里不必假设每个有限 \(p\) 的范数均有限；允许无穷值正好覆盖本性无界函数。
\end{solution}

\begin{exercise}\label{b1:ex:10-interpolation}
设 \(1\leq p<r\leq\infty\)、\(0<\theta<1\)，并由
\(1/q=\theta/p+(1-\theta)/r\) 定义 \(q\)。证明对任意测度空间上的 \(f\in L^p\cap L^r\)，有
\[
 \|f\|_q\leq\|f\|_p^\theta\cdot\|f\|_r^{1-\theta}.
\]
\end{exercise}
\begin{solution}
若 \(r<\infty\)，两个数 \(p/(\theta q)\)、\(r/\bigl((1-\theta)q\bigr)\) 互为共轭指数。
将 Hölder 不等式用于 \(|f|^{\theta q}\) 与 \(|f|^{(1-\theta)q}\)，得到
\[
 \int|f|^q\,d\mu\leq
 \left(\int|f|^p\,d\mu\right)^{\theta q/p}
 \left(\int|f|^r\,d\mu\right)^{(1-\theta)q/r}.
\]
取 \(q\) 次方根即可。若 \(r=\infty\)，则 \(\theta=p/q\)，且
\(\int|f|^q\leq\|f\|_\infty^{q-p}\int|f|^p\)，仍得所求不等式。
该估计不要求测度有限：同时拥有两端的积分控制，足以得到中间指数的控制。
\end{solution}

\begin{exercise}[对偶范数]\label{b1:ex:10-dual-norm-test}
设 \(\mu\) 为 \(\sigma\)-有限正测度，\(1\leq p<\infty\)、\(q=p'\)，且 \(g\in L^q(\mu;\mathbb K)\)。定义
\[
 \Lambda_g(f)=\int_Xfg\,\mathrm d\mu,
 \qquad f\in L^p(\mu;\mathbb K).
\]
不用对偶表示定理的满射部分，直接证明 \(\Lambda_g\in(L^p)^*\) 且
\[
 \|\Lambda_g\|=\|g\|_q.
\]
当 \(1<p<\infty\) 时找出达到该范数的单位测试函数；当 \(p=1\) 时构造一族使泛函值趋近范数的单位测试函数，并说明 \(\sigma\)-有限性在何处使用。
\end{exercise}
\begin{solution}
Hölder 不等式先给出 \(\|\Lambda_g\|\leq\|g\|_q\)。若 \(g=0\)，结论已经成立。以下设 \(g\ne0\)，并令
\[
 \vartheta_g=
 \begin{cases}
  \overline g/|g|,&g\ne0,\\
  0,&g=0.
 \end{cases}
\]
当 \(1<p<\infty\) 时，取
\[
 f_0=\frac{\vartheta_g|g|^{q-1}}{\|g\|_q^{q-1}}.
\]
由 \((q-1)p=q\) 可知 \(\|f_0\|_p=1\)，而
\(\Lambda_g(f_0)=\|g\|_q\)，所以范数在 \(f_0\) 处达到。

当 \(p=1\)、\(q=\infty\) 时，任取 \(0<t<\|g\|_\infty\)。集合 \(B_t=\{|g|>t\}\) 有正测度。把 \(X\) 写成可数个有限测度集的并后，至少有一个交集给出可测集 \(A_t\subset B_t\)，满足
\(0<\mu(A_t)<\infty\)。令
\[
 f_t=\frac{\vartheta_g\mathbf1_{A_t}}{\mu(A_t)}.
\]
则 \(\|f_t\|_1=1\)，且
\[
 \Lambda_g(f_t)=\frac1{\mu(A_t)}\int_{A_t}|g|\,\mathrm d\mu>t.
\]
令 \(t\uparrow\|g\|_\infty\)，得到反向范数估计。这里 \(\sigma\)-有限性只用于从正测度集合 \(B_t\) 中找到正而有限测度的测试集；有限指数情形的显式测试不需要这一假设。
\end{solution}

\begin{exercise}\label{b1:ex:10-measure-to-weak}
设 \(\mu(\Omega)<\infty\)、\(1<p<\infty\)，且 \(\sup_n\|f_n\|_p<\infty\)。
若 \(f_n\) 依测度趋于有限值可测函数 \(f\)，证明 \(f\in L^p\) 且 \(f_n\rightharpoonup f\) 于 \(L^p\)。
说明把 \(p>1\) 改成 \(p=1\) 后结论可以失败。
\end{exercise}
\begin{solution}
由\cref{b1:lem:measure-convergence-subsequence}可抽取几乎处处收敛子列，
Fatou 引理给出 \(\|f\|_p\leq M=\sup_n\|f_n\|_p\)。
对任意可测 \(E\)，Hölder 不等式给出
\[
 \int_E|f_n|\,d\mu\leq M\mu(E)^{1-1/p}.
\]
由于测度有限，\(L^1\) 范数也一致有界，所以 \((f_n)\) 一致可积。
由\cref{b1:thm:vitali-convergence}，\(\|f_n-f\|_1\to0\)。

令 \(q=p/(p-1)\)。对任意 \(g\in L^q\)，取有界截断 \(g_k=g\mathbf1_{\{|g|\leq k\}}\)，则
\[
 \left|\int(f_n-f)g\,d\mu\right|
 \leq k\|f_n-f\|_1+2M\|g-g_k\|_q.
\]
先固定 \(k\) 令 \(n\to\infty\)，再令 \(k\to\infty\)，得到积分测试收敛。
对偶表示证明这就是弱收敛。
在 \((0,1)\) 上取 \(f_n=n\mathbf1_{(0,1/n)}\)，则 \(f_n\to0\) 依测度且 \(\|f_n\|_1=1\)，
但用常数函数一测试，积分始终为一，故不弱趋于零。
\end{solution}

\begin{exercise}\label{b1:ex:10-multiplication}
设 \(0<\mu(\Omega)<\infty\)、\(a\in L^\infty(\mu)\)。证明乘法算子
\(M_a:L^p\to L^p\)、\(M_af=af\) 对每个 \(1\leq p\leq\infty\) 都有界，且 \(\|M_a\|=\|a\|_\infty\)。
\end{exercise}
\begin{solution}
由几乎处处估计 \(|af|\leq\|a\|_\infty\cdot|f|\)，得 \(\|M_a\|\leq\|a\|_\infty\)。
若 \(0<c<\|a\|_\infty\)，则 \(E=\{|a|>c\}\) 有正的有限测度。
在有限 \(p\) 时用 \(f=\mu(E)^{-1/p}\mathbf1_E\) 测试，有 \(\|f\|_p=1\) 且 \(\|af\|_p\geq c\)。
在 \(p=\infty\) 时则用 \(f=\mathbf1_E\)，同样得到下界。
令 \(c\uparrow\|a\|_\infty\) 即得等号；\(a=0\) 的情形直接成立。
有限测度保证这些集合的特征函数都是合法测试，而不只是形式上的试代。
\end{solution}

\begin{exercise}\label{b1:ex:10-translation-continuity}
设 \(1\leq p<\infty\)、\(f\in L^p(\mathbb R^d)\)，令 \(\tau_hf(x)=f(x-h)\)。
证明 \(\|\tau_hf-f\|_p\to0\) 当 \(h\to0\)。举例说明这在 \(L^\infty\) 中一般不成立。
\end{exercise}
\begin{solution}
先设 \(g\in C_c(\mathbb R^d)\)。当 \(|h|\leq1\) 时，\(g\) 与 \(\tau_hg\) 的支集落在同一个有界集合中；
由一致连续性，\(\|\tau_hg-g\|_\infty\to0\)，乘上该集合测度的 \(1/p\) 次方，得 \(L^p\) 收敛。
对一般 \(f\)，选 \(g\in C_c\) 使 \(\|f-g\|_p<\varepsilon\)。
Lebesgue 测度的平移不变性给出 \(\|\tau_h(f-g)\|_p=\|f-g\|_p\)，因而
\[
 \|\tau_hf-f\|_p\leq2\varepsilon+\|\tau_hg-g\|_p.
\]
先令 \(h\to0\)，再令 \(\varepsilon\downarrow0\)。
另一方面，对 \(f=\mathbf1_{(0,1)}\)，任意 \(0<|h|<1\) 都有 \(\|\tau_hf-f\|_\infty=1\)。
跳跃附近的集合虽越来越小，本性上确界仍看得见完整的跳跃高度。
\end{solution}

\begin{exercise}\label{b1:ex:10-minimal-norm}
设 \(\mu\) 为 \(\sigma\)-有限测度、\(1<p<\infty\)，且 \(C\subset L^p(\mu)\) 是非空范数闭凸集。
证明 \(C\) 中存在唯一的最小范数元素。给出 \(\ell^1\) 和 \(\ell^\infty\) 中唯一性失败的有限维例子。
\end{exercise}
\begin{solution}
选 \(f_n\in C\) 使 \(\|f_n\|_p\to\inf_C\|f\|_p\)。该列有界，自反性给出弱收敛子列。
闭凸集弱闭，故极限 \(f\in C\)；范数弱下半连续，故 \(f\) 取得下确界。

若 \(u\ne v\) 于 \(L^p\)，则 \(u,v\) 在一个正测度集合上不同。
对不同的实数或复数 \(a,b\)，由标量三角不等式及 \(t\mapsto t^p\) 的严格凸性，有
\[
 \left|\frac{a+b}{2}\right|^p<\frac{|a|^p+|b|^p}{2}.
\]
若模不同，严格性来自后者；若模相同而数不同，来自前者。
逐点积分可知 \(\|(u+v)/2\|_p^p<(\|u\|_p^p+\|v\|_p^p)/2\)。
两个不同极小元的中点因而会有更小范数，矛盾。

在 \(\mathbb R^2\) 的 \(\ell^1\) 范数下，连接 \((1,0)\) 与 \((0,1)\) 的整个线段范数恒为一。
在 \(\ell^\infty\) 范数下，连接 \((1,-1)\) 与 \((1,1)\) 的线段亦如此。
将它们嵌入相应序列空间，就得到闭凸约束下极小元不唯一的例子。
\end{solution}

\begin{optionalexercise}[集中、逃逸与振荡]\label{b1:ex:10-endpoint-three-mechanisms}
分别考虑
\[
 c_n=n\mathbf1_{(0,1/n)}\quad\text{于 }(0,1),
 \qquad
 e_n=\mathbf1_{[n,n+1)}\quad\text{于 }\mathbb R,
\]
以及 \(r_n(x)=(-1)^{\lfloor2^nx\rfloor}\) 于 \((0,1)\)。

（1）验证三列都在 \(L^1\) 中有界，并判断它们是否满足本书定义的幅值尾部一致可积性。

（2）利用 Dunford--Pettis 定理和有界测试函数，说明 \((c_n)\)、\((e_n)\)、\((r_n)\) 分别表现为质量集中、向无穷远逃逸与振荡；具体判断各列是否有弱收敛子列，并说明 \((e_n)\) 还缺少哪一项空间控制。

（3）对集中列取 \(E_k=(0,1/k)\)，验证它们可作为 Chacon 引理中的递减坏集，并指出坏集外的极限为何没有保留原列的总积分。
\end{optionalexercise}
\begin{solution}
三列的 \(L^1\) 范数都等于一。若 \(M>1\)，则 \(\{|e_n|>M\}\) 与 \(\{|r_n|>M\}\) 均为空，所以后两列满足幅值尾部条件。对任意 \(M>0\)，选 \(n>M\)，便有
\[
 \int_{\{|c_n|>M\}}|c_n|\,\mathrm dx=1,
\]
所以集中列不一致可积。

在有限测度空间 \((0,1)\) 上，Dunford--Pettis 定理立即说明 \(\{c_n:n\geq1\}\) 不相对弱紧。还可直接看出它没有弱收敛子列：若某子列弱收敛于 \(c\)，则对每个 \(\delta>0\) 以及支集包含在 \((\delta,1)\) 中的有界测试函数，积分最终为零，故 \(c=0\) 几乎处处；但用常函数一测试又会得到 \(\int_0^1c=1\)，矛盾。

函数 \(r_n\) 在 \(L^2(0,1)\) 中构成标准正交系。对每个 \(h\in L^\infty(0,1)\subset L^2(0,1)\)，Bessel 不等式给出 \(\int_0^1r_nh\,\mathrm dx\to0\)；复值 \(h\) 时对 \(\overline h\) 使用同一论证。因此 \(r_n\rightharpoonup0\) 于 \(L^1\)，集合 \(\{r_n:n\geq1\}\cup\{0\}\) 是弱紧的。另一方面，\(m\ne n\) 时 \(\|r_m-r_n\|_1=1\)，所以这只是弱收敛，不是范数紧性。

逃逸列虽满足幅值尾部条件，却不满足\eqref{b1:eq:l1-no-escape}。事实上，对每个有限测度集 \(E\subset\mathbb R\)，半开区间两两不交给出
\[
 \sum_{n=1}^\infty m\bigl(E\cap[n,n+1)\bigr)\leq m(E)<\infty,
\]
故这些交集的测度趋于零，而 \(\int_{\mathbb R\setminus E}|e_n|\to1\)。若某个子列弱收敛于 \(e\)，用任意有界支集的 \(L^\infty\) 函数测试可知 \(e=0\) 几乎处处；再用常函数一测试则得到 \(0=1\)，矛盾。因此它也没有弱收敛子列，缺少的正是空间尾部控制。

最后，\(E_k\downarrow\varnothing\) 且 \(m(E_k)=1/k\to0\)。固定 \(k\) 后，只要 \(n\geq k\)，便有 \(c_n\mathbf1_{(0,1)\setminus E_k}=0\)，所以坏集外的整族只含有限多个非零函数，因而一致可积并弱趋于零。然而
\(\int_{E_k}c_n=1\) 对所有 \(n\geq k\) 成立；全部质量都留在被删去的坏集内，故局部化弱极限不保留原列恒为一的总积分。
\end{solution}

\begin{exercise}\label{b1:ex:10-separability}
证明 \(L^p(\mathbb R^d)\) 在 \(1\leq p<\infty\) 时可分，而 \(L^\infty(0,1)\) 不可分。
\end{exercise}
\begin{solution}
取有限多个端点坐标为有理数的有界半开长方体的特征函数，以有理系数作线性组合；
复情形允许实部、虚部均为有理数的系数。所得函数族可数。
它在 \(C_c\) 中关于 \(L^p\) 范数稠密：把给定连续紧支集函数的支集放在一个有理长方体内，
用足够细的有理网格分割它，按一致连续性在每格选近似常值，再把这些常值换成有理系数。
网格边界为零测集，统一误差乘以固定长方体测度的 \(1/p\) 次方即给出 \(L^p\) 误差。
结合 \(C_c\) 的稠密性，得到全空间可分。

反之，令 \(u_t=\mathbf1_{(0,t)}\)、\(0<t<1\)。当 \(s\ne t\) 时，\(\|u_s-u_t\|_\infty=1\)。
任意半径 \(1/3\) 的球至多包含一个这样的函数。
若存在可数稠密集，以其成员为中心的半径 \(1/3\) 的球应覆盖全空间，却不能覆盖不可数族 \(\{u_t\}\)，矛盾。
这也说明本性上确界拓扑中不能照搬有限 \(p\) 的可数逼近方案。
\end{solution}
